\documentclass[10pt,a4paper]{amsart}
\usepackage[margin=19mm]{geometry}
\usepackage{amsmath,amssymb,mathtools}
\usepackage[scr=rsfs,cal=euler]{mathalfa}
\usepackage{xfrac}
\usepackage[unicode,hidelinks,bookmarks=false]{hyperref}
\newcommand{\ie}{i.e.\ }
\newcommand{\cf}{cf.\ }
\newcommand{\resp}{respectively\ }
\newcommand{\Subsection}{\subsection}
\providecommand{\nobreakdash}{\nobreak\hbox{-}}
\setlength{\emergencystretch}{2em}
\multlinegap0pt
\usepackage{xcolor}
\usepackage{amssymb}
\usepackage{enumerate}
\usepackage[all]{xy}
\usepackage{mathtools}
\usepackage[shortlabels]{enumitem}
\usepackage{bm}
\usepackage{bbm}
\usepackage[normalem]{ulem}
\usepackage{tikz}
\newtheorem{lemma}{Lemma}
\def\N{{\mathbb N}}
\newcommand{\kl}{\left(}
\newcommand{\kr}{\right)}
\newcommand{\tm}{\subseteq}
\title{Constructing Toeplitz arrays via cut and project schemes}
\author{Mar{\'i}a Isabel Cortez, Jamal Drewlo, Jaime G{\'o}mez, Tobias J{\"a}ger}
\date{Source: arXiv:2512.19921v1 (CC BY 4.0)}
\begin{document}
\maketitle

\noindent\emph{Source and licence.} Constructing Toeplitz arrays via cut and project schemes. Version arXiv:2512.19921v1 (CC BY 4.0), distributed by arXiv under the Creative Commons Attribution 4.0 International licence, http://creativecommons.org/licenses/by/4.0/, as recorded in the arXiv licence metadata for this exact version and shown on the abstract page https://arxiv.org/abs/2512.19921v1. Full licence notice: \texttt{licenses/arxiv-2512.19921-CC-BY-4.0.txt}.\par\smallskip

\noindent\textit{Editorial scope.} Counted unit: the article's lemma Borel\_map (source0.tex lines 1624--1627). The article's complete original proof of it is delivered with it (source0.tex lines 1628--1662, one \texttt{proof} environment of 35 lines). No mathematics is written, repaired or re-derived: the statement and the proof are the article's own lines, in the article's own notation, and no retained line is substituted or re-flowed.  Retained uncounted as the source-local context the proof needs: lines 1205--1208; lines 1210--1215; lines 1260--1274.  Nothing was omitted from any retained span, so the package carries no omission record.  No retained line contains a citation, so the wrapper carries no bibliography and no bibliography entry.  No retained line contains a figure environment, an image-inclusion command or drawing code, so no image is delivered and the package declares no asset dependency.  Editorial formatting only, in addition: an AMS \texttt{amsart} preamble that restates the article's own declarations and macros used by the retained lines, namely the environments \texttt{lemma} on separate counters, together with the standard packages those lines need.  The retained environments print in this document as Lemma 1 in order of appearance, whereas the article prints the same items with the numbering of its own full document.  The complete native source of the article as distributed in the arXiv e-print for this exact version is bundled unchanged as \texttt{sources/191520/source0.tex}.
\begin{equation}
    \label{eq: definition W}
    W = \overline{U}.
\end{equation}

\begin{lemma}\label{lem: alternative chara W and partial W}
The set $W\subseteq \overleftarrow{G}$ defined in (\ref{eq: definition W}) is proper, with $\mathrm{int}(W)=U$ and 
 \begin{equation}\label{eq: boundary characterisation}
 \partial W=\bigcap_{n\in\mathbb{N}}Z_n=\overleftarrow{G}\setminus (V\cup U).
 \end{equation}
\end{lemma}

\begin{lemma}\label{lem: chara nonempty intersect with W and complement}
    Let $C\tm \overleftarrow{G}$ be a cylinder of level $n$.
    The following statements are equivalent:
    \begin{enumerate}
        \item $C\cap W\neq \emptyset \neq  C\cap (\overleftarrow{G}\setminus W)$.
        \item $C\cap Z_n \neq \emptyset$.
        \item $C\tm Z_n$.
        \item $C\cap \partial W\neq \emptyset$.
    \end{enumerate}
    Furthermore, the following statements are equivalent:
    \begin{enumerate}
        \item $\#\mathcal{M}_n(W,C)=1$ for every cylinder $C$ of level $n$ with $C\cap W\neq \emptyset \neq  C\cap (\overleftarrow{G}\setminus W)$.
        \item $\# B_{n+1} =1$.
    \end{enumerate}
\end{lemma}

\begin{lemma}\label{Borel_map}
For every $1\leq j\leq k+1$, the map $\phi_j:H_0\to \overline{O_G(x_{W^{(k)}})}=X$ defined above 
is Borel.
\end{lemma}

\begin{proof}
Let $n\in\N$ and $a_{n,1},\ldots, a_{n,k_n}\in G$ be such that $[a_{n,1}]_n,\ldots, [a_{n,k_n}]_{n}$  are  the different cylinder sets of level $n$ (where $k_n=[G:\Gamma_n]$). For every $1\leq j\leq k+1$, define 
$$W_{n,j}^{(k)}=\kl\mathrm{int}(W^{(k)})\setminus\bigcup_{i=1}^{j-1} (H_i \cap Z_n)\kr\cup \bigcup_{i=j}^k\left(H_i\cap Z_n\right),$$
and for every $\xi\in H_0$ and $g\in G$ we set $$\phi_{n,j}(\xi) = x_{W_{n,j}^{(k)}\psi_n(\xi)^{-1}},$$ which means that 
$$\phi_{n,j}(\xi)(g) = 1 \iff \tau(g) \in W_{n,j}^{(k)}\psi_n(\xi)^{-1}.$$

Since  $\phi_{n,j}(H_0)=\left\{x_{W_{n,j}^{(k)}a_{n,s}^{-1}}: 1\leq s\leq k_n\right\}$ and
$$(\phi_{n,j})^{-1}\left(\left\{x_{W_{n,j}^{(k)}a_{n,s}^{-1}}\right\}\right)=[a_{n,s}]_n,$$ for every $1\leq s\leq k_n$ and $n\in \N$, the maps $\phi_{n,j}$ are Borel.
It is therefore enough to show that $\lim_{n\to \infty}\phi_{n,j}(\xi)=\phi_j(\xi)$, for every $\xi\in H_0$. 

Let $\xi\in H_0$ and $g\in G$.
There are four cases to distinguish:
\begin{enumerate}
    \item $\phi_j(\xi)(g)=1$ and $\tau(g)\xi \in \partial W^{(k)}$.
    \item $\phi_j(\xi)(g)= 1 $ and $\tau(g) \xi\notin \partial W^{(k)}$.
    \item $\phi_j(\xi)(g)=0$ and $\tau(g)\xi \in \partial W^{(k)}$. 
    \item $\phi_j(\xi)(g)=0$ and $\tau(g)\xi\notin \partial W^{(k)}$.
\end{enumerate}
In case (1), we see that $\tau(g)\xi \in \bigcup_{i=j}^k(\partial W^{(k)} \cap H_i)$.
Let $j\leq i \leq k$ such that $\tau(g)\xi \in \partial W^{(k)}\cap H_i$.
Due to Lemma~\ref{lem: alternative chara W and partial W}, this implies for every $n\geq L$ that $\tau(g)\xi\in Z_n \cap H_i$.
Now, we have $[\tau(g)\xi]_n = [\tau(g)\psi_n(\xi)]_n$.
Since $Z_n\cap H_i$ is a union of cylinder sets of level $n$, this also implies $\tau(g) \psi_n(\xi) \in Z_n\cap H_i \tm W_{n,j}^{(k)}$.
Therefore, we have $\phi_{n,j}(\xi)(g) = 1$ for all $n\geq L$, which shows that $\lim_{n\to \infty}\phi_{n,j}(\xi)(g)=\phi_j(\xi)(g)$ in this case.\\
In case (2) we obtain $\tau(g)\xi \in \mathrm{int}(W^{(k)})$.
Choose $m\in \N$ such that $[\tau(g)\xi]_m\tm \mathrm{int}(W^{(k)})$.
In particular, we have $[\tau(g)\xi]_m\cap \partial W_{\mathrm{perf}} = [\tau(g)\xi]_m \cap \partial W^{(k)} = \emptyset$.
Thus, Lemma~\ref{lem: chara nonempty intersect with W and complement} yields $[\tau(g)\xi]_m\cap Z_m=\emptyset$.
Altogether, we obtain for all $n\geq m$ that 
$$ \tau(g)\psi_n(\xi) \in [\tau(g)\xi]_m \tm \mathrm{int}(W^{(k)}) \setminus Z_m \tm W_{n,j}^{(k)},$$
so that $\phi_{n,j}(\xi)(g) = 1$. This again shows $\lim_{n\to \infty}\phi_{n,j}(\xi)(g)=\phi_j(\xi)(g)$.\\
In case (3) we deduce $\tau(g)\xi\in \bigcup_{i=1}^{j-1}(\partial W^{(k)}\cap H_i)$ and can proceed in an analougous way as in (1) to show $\phi_{n,j}(\xi)(g) = 0$ for all $n\geq L$. Finally, in case (4) we obtain $\tau(g)\xi \in \overleftarrow{G}\setminus W^{(k)}$ which is analogous to (2).

In conclusion, we have shown that $\phi_j:H_0\to X$ is a pointwise limit of Borel maps, so that $\phi_j$ itself is Borel.
\end{proof}

\end{document}
